\documentclass[11pt]{article}
\usepackage[margin=1in]{geometry}
\usepackage{amsmath,amssymb,amsthm,booktabs}
\usepackage[hidelinks]{hyperref}

\newtheorem{theorem}{Theorem}
\newtheorem{lemma}[theorem]{Lemma}
\newtheorem{proposition}[theorem]{Proposition}
\newcommand{\R}{\mathbb R}
\newcommand{\re}{\operatorname{Re}}
\newcommand{\impart}{\operatorname{Im}}
\newcommand{\ii}{\mathrm i}

\title{A scale-free proof of the connector seam sign at Poincar\'e's point \(D\)}
\author{}
\date{August 2026}

\begin{document}
\maketitle

\begin{abstract}
The last unresolved numerical issue in the current connector construction is caused by a
double zero at the collision point \(D\). Direct interval evaluation loses the small common
Morse scale and produces boxes containing zero. We retain that scale algebraically. The
oriented real derivative is written exactly as
\[
  L\,\re X_s(L,\delta)+\delta^2\,\re Z_s(L,\delta).
\]
Here \(s=-1\) and \(s=+1\) denote the initial and final connectors. An independent
outward-rounded Python interval calculation proves
\(\re X_s\geq0.00419\) and \(\re Z_s\geq0.06558\) on both complete collars.
Since \(L>0\), the required derivative orientation follows uniformly, however small the
nonconstructively selected Morse length is. Ordinary one-variable calculus then supplies
the missing connector seam sign.
\end{abstract}

\section{Scope}

This note addresses the local-to-global seam problem for the two affine connectors at the
specific collision point \(D\). It uses analytic facts already established in the project:
the inverse-Morse endpoint can be chosen with the normalized bounds in
\S\ref{sec:boxes}, the companion collision factor is positive on the collars, and the first
factor has the required imaginary monotonicity. The new content is the real derivative sign
for the first factor.

This does not prove the global source-contour deformation in \S\S95--100 or the
source-specific boundary-logarithm reduction in \S102. It removes the scale-sensitive seam
obstruction that occurs before those broader tasks.

\section{The centered collision factor}

Let \(x_D\) be the unique root in \([-0.027,-0.026]\) of
\[
  2500x^3+500025x^2+12501x-25=0,
\]
let \(D<0\) be its real cube root, and put
\[
  y_D=\frac{(x_D-1/100)^2}{2(1+(1/100)^2)x_D}.
\]
Write \(F(\zeta,u)\) for the collision factor which vanishes twice at
\((\zeta_D,D)\), and write \(G=\partial F/\partial u\). Direct algebra, with every
inverse-power difference cancelled before evaluation, gives
\begin{equation}\label{eq:centered-expansion}
 G(\zeta,u)
 =A(u-D)+(u-D)^2K(u)
   +\left(\frac{\zeta}{\zeta_D}-1\right)P(u),
\end{equation}
where
\[
 A=\frac{\partial^2F}{\partial u^2}(\zeta_D,D)\in\R,
 \qquad A<0.
\]
The explicit formulas used by the calculation are in Appendix~\ref{sec:formulas}.

Set
\[
 E(a)=
 \begin{cases}
  0,&a=0,\\
  \dfrac{e^a-1-a}{a^2},&a\ne0.
 \end{cases}
\]
Then
\begin{equation}\label{eq:exp-remainder}
 e^a=1+a+a^2E(a),
 \qquad |a|\leq1\Longrightarrow |E(a)|\leq1.
\end{equation}
The nonstandard value at \(a=0\) is harmless because it is multiplied by \(a^2\).

\section{Scale-free connector coordinates}\label{sec:boxes}

Let \(s=-1\) on the initial connector and \(s=+1\) on the final connector. Let \(L>0\)
be the selected Morse half-length and let \(\delta\) be distance from the local endpoint:
\[
 \delta=1-t\quad(s=-1),\qquad \delta=t\quad(s=+1).
\]
On both collars, \(0\leq\delta\leq261/1024\). Normalize the moving inputs by
\[
 u_{\mathrm{loc},s}-D=sLq_s,\qquad
 u_{\mathrm{out},s}(\kappa)-u_{\mathrm{out},s}(0)=L^2r_s,\qquad
 \frac{\zeta}{\zeta_D}-1=L^2p.
\]
The inverse-Morse model can be shrunk so that
\begin{equation}\label{eq:primitive-bounds}
 0<L\leq2^{-20},\qquad
 |\re r_s|,|\impart r_s|,|\re p|,|\impart p|\leq1,
\end{equation}
and
\begin{equation}\label{eq:q-box}
 -\frac{1024}{2^{20}}\leq\re q_s\leq\frac{1024}{2^{20}},
 \qquad
 -\frac{158311}{2^{20}}\leq\impart q_s\leq-\frac{130048}{2^{20}}.
\end{equation}
This last box is the quantitative form of the downward inverse-Morse tangent.

At collapsed scale the local-to-outer directions are
\[
 B_-=-D(1+\ii),\qquad B_+=-D(1-\ii).
\]
Define the actual direction and coordinate by
\begin{equation}\label{eq:direction-coordinate}
 V_s=B_s+L^2r_s-sLq_s,\qquad
 u-D=sLq_s+\delta V_s.
\end{equation}
These are identities. In particular, all occurrences of the small endpoint displacement
share the same \(L\).

\section{The homogeneous real-part identity}

Substitute \eqref{eq:direction-coordinate} and the normalized parameter displacement into
\eqref{eq:centered-expansion}. Define
\begin{align}
 X_s={}&Asq_sV_s
 +\delta A(Lr_s-sq_s)(V_s+B_s) \notag\\
 &+\bigl(Lq_s^2+2s\delta q_sV_s\bigr)K(u)V_s
 +LpP(u)V_s,\label{eq:X}\\[3pt]
 Z_s={}&V_s^3K(u).\label{eq:Z}
\end{align}

\begin{lemma}[Exact cancellation]\label{lem:cancellation}
For either connector,
\begin{equation}\label{eq:homogeneous-identity}
 \re\bigl(G(\zeta,u)V_s\bigr)=L\re X_s+\delta^2\re Z_s.
\end{equation}
\end{lemma}

\begin{proof}
Expand the left side using \eqref{eq:centered-expansion} and
\(u-D=sLq_s+\delta V_s\). Also write
\(V_s=B_s+L(Lr_s-sq_s)\). Collect every term containing \(L\) into
\(LX_s\), and collect the quadratic-distance term into \(\delta^2Z_s\).
The only apparent term left over is \(\delta AB_s^2\). Since \(A\) is real and
\[
 B_-^2=2\ii D^2,\qquad B_+^2=-2\ii D^2,
\]
its real part is exactly zero.
\end{proof}

Neither coefficient in \eqref{eq:homogeneous-identity} is divided by \(L\). The first term
retains a positive factor \(L\), while the collapsed part begins at order \(\delta^2\).
Thus the coefficient boxes can remain separated from zero as \(L\to0^+\).

\section{Finite interval lemma}

The primitive dyadic boxes, all at denominator \(2^{20}\), are
\begin{align*}
 D&\in[-314053,-314047]2^{-20},&
 y_D&\in[-28312,-25000]2^{-20},\\
 \frac{100}{30003}&\in[3494,3495]2^{-20},&
 \frac{200}{10001}&\in[20969,20970]2^{-20},\\
 \frac1{10001}&\in[104,105]2^{-20}.&&
\end{align*}
The cubic is strictly decreasing on
\([-0.026865395705,-0.026865395704]\), takes opposite signs at its endpoints,
and its negative real cube root lies in the displayed \(D\)-box. Substitution in the
formula for \(y_D\) gives the displayed \(y_D\)-box.

Cover the distance interval by the 160 uniform cells
\[
 I_j=[1671j/2^{20},1671(j+1)/2^{20}],\qquad 0\leq j<160.
\]
Indeed, (160\cdot1671/2^{20}=0.254974365234375>261/1024).
On each cell and side, rectangular complex interval arithmetic evaluates
\eqref{eq:direction-coordinate}, the Appendix formulas, and \eqref{eq:X}--\eqref{eq:Z}.
The exponential remainder \(E(a)\) is enclosed by
\([-1,1]+\ii[-1,1]\), using \eqref{eq:exp-remainder}; the same interval evaluation first
checks \(|a|\leq0.218249762481<1\) on every cell. All other normalized moving
inputs use \eqref{eq:primitive-bounds}--\eqref{eq:q-box} independently, deliberately
overestimating their joint range.

\begin{proposition}[Interval result]\label{prop:interval}
For \(s=-1,+1\), every admissible normalized input, and
\(0\leq\delta\leq261/1024\),
\[
 \re X_s\geq0.0041922727090,\qquad
 \re Z_s\geq0.0655826074416.
\]
\end{proposition}

The independently computed enclosures are:
\begin{center}
\begin{tabular}{lcc}
\toprule
 & initial connector & final connector\\
\midrule
minimum lower endpoint of \(\re X_s\) & 0.00419227270907 & 0.00419227270907\\
maximum upper endpoint of \(\re X_s\) & 0.218433271265 & 0.218433271265\\
minimum lower endpoint of \(\re Z_s\) & 0.0655826074417 & 0.0655826074417\\
maximum upper endpoint of \(\re Z_s\) & 0.776610992074 & 0.776610992074\\
\bottomrule
\end{tabular}
\end{center}
The worst lower bound for \(X_s\) occurs in cell \(159\), at the outer collar boundary;
the worst lower bound for \(Z_s\) occurs in cell \(0\), at the local endpoint.
The two row lists agree exactly after conjugating the imaginary coordinate, so the checked
one-side table in \texttt{research/chapter\_vi\_connector\_seam\_table.csv} represents both
connectors. The verifier recomputes this symmetry rather than assuming it.

The calculation is reproduced by
\begin{verbatim}
python3 research/check_chapter_vi_connector_seam.py
\end{verbatim}
The script uses only the Python standard library. Every elementary binary64 result is widened
by one representable number toward each infinity. This is strong independent evidence and is
close to a conventional interval proof, but it is not a kernel-checked certificate.

\section{Conclusion for the seam}

\begin{theorem}[Oriented derivative on both collars]\label{thm:orientation}
For a sufficiently small anchored Morse model satisfying
\eqref{eq:primitive-bounds}--\eqref{eq:q-box},
\[
 -\frac{d}{dt}\re F(\zeta,u_-(t))>0
 \quad\text{on the initial collar},
\]
and
\[
 \frac{d}{dt}\re F(\zeta,u_+(t))>0
 \quad\text{on the final collar}.
\]
\end{theorem}

\begin{proof}
The oriented derivative is the left side of \eqref{eq:homogeneous-identity}.
Proposition~\ref{prop:interval}, \(L>0\), and \(\delta^2\geq0\) give
\[
 L\re X_s+\delta^2\re Z_s
 \geq0.0041922727090L+0.0655826074416\delta^2>0.
\]
The orientation reverses the initial connector and leaves the final connector unchanged.
\end{proof}

At the local endpoint, the real part of the first factor is already strictly positive:
the product radicand is positive real there and the companion factor has positive real part.
On the initial connector, Theorem~\ref{thm:orientation} says that the first factor's real part
is nonincreasing as \(t\) approaches \(1\), so every earlier collar value is at least its
positive endpoint value. On the final connector it is nondecreasing away from \(t=0\).
Consequently
\[
 \re F(\zeta,u_s(t))>0
 \quad\text{throughout both endpoint collars}.
\]

In particular, if the imaginary part of the first factor vanishes at its unique possible
crossing, its real part is positive. The principal square root therefore extends from the
outer boundary to the prescribed local Morse root without crossing its branch cut. The
existing positive-real result for the companion factor supplies the other square root.
Their product gives the compatible radicand sheet on each connector, and the two local seam
signs are the required positive Morse signs. Together with the already constructed outer arcs
and pinched middle arc, this supplies the seam-compatible five-piece logarithmic contribution.

\appendix
\section{Explicit coefficient formulas}\label{sec:formulas}

Put
\[
 c=\frac{100}{30003},\qquad b=\frac{200}{10001},\qquad
 Q=u^2+uD+D^2,\qquad Q_1=u+2D.
\]
The coefficient of \(u-D\) in the relative exponential argument and the multiplier are
\begin{align*}
 H(u)&=-cQ\bigl(1+u^{-3}D^{-3}\bigr),\\
 M(u)&=-2y_DD^{-1}+by_DD^{-1}(u^{-3}+u^3).
\end{align*}
The argument in \eqref{eq:exp-remainder} is \(a=(u-D)H(u)\), and
\[
 P(u)=\bigl(1+a+a^2E(a)\bigr)M(u).
\]

For the quadratic coefficient, introduce
\begin{align*}
 \eta_2&=-\frac{u+D}{u^2D^2},&
 \eta_4&=-\frac{u^3+u^2D+uD^2+D^3}{u^4D^4},\\
 T&=D^{-4}(u^{-2}+D^2u^{-4}),&
 T_1&=D^{-4}(\eta_2+D^2\eta_4),\\
 R_L(u)&=\frac1{10001}\bigl((30000+3T)+6DT_1\bigr),\\
 \eta_3&=-\frac{Q}{u^3D^3},\\
 R_S(u)&=D^{-1}\left[Q_1(1-u^{-3}D^{-3})
      +3D^2(-D^{-3}\eta_3)\right],\\
 H_1(u)&=-c\left[Q_1(1+u^{-3}D^{-3})+3D^2(D^{-3}\eta_3)\right],\\
 M_1(u)&=by_DD^{-1}(\eta_3+Q).
\end{align*}
Then
\[
 K(u)=R_L(u)+by_DR_S(u)+H_1(u)M(u)+H(D)M_1(u)+H(u)^2E(a)M(u).
\]

Finally, define
\begin{align*}
 C_L(u)&=\frac{u+D}{10001}
 \left(30000+3\frac{u^2+D^2}{u^4D^4}\right),\\
 C_S(u)&=QD^{-1}(1-u^{-3}D^{-3}),\\
 C(u)&=C_L(u)+by_DC_S(u).
\end{align*}
The base coefficient is evaluated without numerical differentiation:
\[
 A=C(D)+H(D)M(D),\qquad
 -2.69848482771\leq A\leq-2.56120107040.
\]

\end{document}
